\begin{table*}[!ht]

\renewcommand{\arraystretch}{1.1}

\small
\centering
    \begin{tabular}{p{1.0\textwidth}}
    \toprule
    \textbf{\textsc{Summarization}}\\
    \midrule
    Below is the original news:\\
\{article\}\\
Below is a summary of the news:\\
\{summary\}\\
Your task is to determine whether the summary contains either or both of the following two types of hallucinations:\\
1. conflict: instances where the summary presents direct contraction or opposition to the original news;\\
2. baseless info: instances where the generated summary includes information which is not substantiated by or inferred from the original news. \\
Then, compile the labeled hallucinated spans into a JSON dict, with a key "hallucination list" and its value is a list of hallucinated spans. If there exist potential hallucinations, the output should be in the following JSON format: \{"hallucination list": [hallucination span1, hallucination span2, ...]\}. Otherwise, leave the value as a empty list as following: \{"hallucination list": []\}.\\
Output:\\
    \midrule
    \textbf{\textsc{Question answering}}\\
    \midrule
        Below is a question:\\
\{question\}\\
Below are related passages:\\
\{passages\}\\
Below is an answer:\\
\{answer\}\\
Your task is to determine whether the answer contains either or both of the following two types of hallucinations: \\
1. conflict: instances where the answer presents direct contraction or opposition to the passages;\\
2. baseless info: instances where the answer includes information which is not substantiated by or inferred from the passages.\\
Then, compile the labeled hallucinated spans into a JSON dict, with a key "hallucination list" and its value is a list of hallucinated spans. If there exist potential hallucinations, the output should be in the following JSON format: \{"hallucination list": [hallucination span1, hallucination span2, ...]\}. Otherwise, leave the value as a empty list as following: \{"hallucination list": []\}.\\
Output:\\
    
    \midrule
    \textbf{\textsc{Data-to-text Writing}}\\
    \midrule
Below is a structured data in the JSON format:\\
\{business info\}\\
Below is an overview article written in accordance with the structured data:\\
\{overview\}\\
Your task is to determine whether the overview contains either or both of the following two types of hallucinations: \\
1. conflict: instances where the overview presents direct contraction or opposition to the structured data;\\
2. baseless info: instances where the generated overview includes information which is not substantiated by or inferred from the structured data.\\
In JSON, "null" or "None" represents an unknown value rather than a negation. \\
Then, compile the labeled hallucinated spans into a JSON dict, with a key "hallucination list" and its value is a list of hallucinated spans. If there exist potential hallucinations, the output should be in the following JSON format: \{"hallucination list": [hallucination span1, hallucination span2, ...]\}. Otherwise, leave the value as a empty list as following: \{"hallucination list": []\}.\\
Output:\\
    \bottomrule
    \end{tabular}
    \caption{Prompts for detecting hallucination for the three types of tasks. In the prompt for data-to-text writing, we clarified that null or None in JSON should be treated as unknown rather than a negation.} 
    \label{tab:prompts}
\end{table*}
